\chapter{Simplify before you differentiate}

\begin{orient}
\textbf{What this chapter is for.} Differentiation is the one genuinely mechanical operation in calculus: every elementary function has a derivative, the rules compose without exception, and no ingenuity is required to apply them. It follows that a derivative problem cannot be hard \emph{as a derivative problem}. If a problem is hard, the difficulty was placed in the algebra, and the algebra is what a good solver attacks first. This chapter teaches the opening move: read the expression as an algebraist for ten seconds before reaching for any rule. Log laws split a monster into a sum of independent terms; a trig identity turns a product into a single sinusoid and makes the seventh derivative a one-liner; a conjugate pair reveals that the function was a constant all along; a radical is a power in disguise; a series, a product, or an infinite nesting has a closed form that must be found before anything is differentiated. At the end you should be unable to start a differentiation without first asking what the expression really is.
\end{orient}

\section{Ten seconds of algebra}

Here is the asymmetry that organises the whole chapter. Integration is genuinely hard: elementary functions have non-elementary antiderivatives, and there is no algorithm you can run without thought. Differentiation has none of that structure. The derivative of an elementary function is elementary, the rules are finitely many, and applying them is bookkeeping. Nobody sets a problem whose difficulty is ``can you apply the chain rule three times''. What they set is a problem whose difficulty is ``can you see that this expression is $\tfrac12\sin 4x$'', dressed up so that the mechanical route exists but costs two pages.

So a hard derivative problem is a disguised algebra problem, and the disguise is always one of a short list of moves. The list is short because the setter has to be able to write the disguised form down, which means the disguise is an identity applied backwards: a sum was recombined into a single logarithm, a sinusoid was expanded into a product, a constant was written as a difference of two nonconstant pieces, a rational power was written as a nest of radicals. Undoing the disguise is the same identity applied forwards, and there are perhaps eight of them.

\begin{keynote}[The ten-second scan]
Before writing a single prime, answer four questions about the expression on the page. Where does $x$ actually occur? Is there a wrapper (a logarithm, an inverse trig function) that an identity will remove? Are there radicals, reciprocal trig functions, or a quotient with one term downstairs, all of which are powers in disguise? And is the function given to you as a \emph{process} (a sum, a product, an infinite nesting) rather than as a formula? Four questions, ten seconds, and they decide whether the problem takes two lines or two pages. The cost of skipping them is not a wrong answer; it is the right answer obtained slowly, with five opportunities to drop a sign.
\end{keynote}

The scan has an order, because the moves are not independent. Constancy is checked first, since a constant function makes every other question moot. Wrapper removal comes next, because it typically converts one product-and-quotient monster into a sum of terms that then need their own scan. Normalisation to powers comes third. And the process forms are checked last only because they are the rarest, not because they are the least profitable: they are the ones where the payoff is largest.

\begin{center}
\begin{tikzpicture}[node distance=6mm and 8mm]
\node[dtest] (a) {Does $x$ occur, and does it\\ survive the obvious identities?};
\node[dleaf, right=14mm of a] (b) {$f$ is constant.\\ $f'\equiv 0$. Stop.};
\node[dtest, below=9mm of a] (c) {Is $f$ presented as a $\sum$, a $\prod$,\\ or a self-similar nesting?};
\node[dleaf, right=14mm of c] (d) {Collapse to a closed form,\\ then re-enter at the top};
\node[dtest, below=9mm of c] (e) {A single $\ln$, $\arctan$ or $\arcsin$\\ wrapping a product, quotient,\\ power, or radical?};
\node[dleaf, right=14mm of e] (f) {Log laws or a half-angle\\ identity: the monster\\ becomes a sum};
\node[dnode, below=9mm of e] (g) {Radicals to rational powers, reciprocal trig to base form,\\ sums over one monomial split termwise. Only now\\ choose a rule system.};
\draw[dedge] (a) -- node[dlab,above]{no} (b);
\draw[dedge] (a) -- node[dlab,left]{yes} (c);
\draw[dedge] (c) -- node[dlab,above]{yes} (d);
\draw[dedge] (c) -- node[dlab,left]{no} (e);
\draw[dedge] (e) -- node[dlab,above]{yes} (f);
\draw[dedge] (e) -- node[dlab,left]{no} (g);
\end{tikzpicture}
\end{center}

Read the tree as a filter, not as a classifier. Most expressions fall through every test and land in the bottom box, which is the honest outcome and the subject of the next chapter. The tree earns its keep on the minority that do not, and on those the saving is not marginal: it is the difference between a line and a page.

\section{Removing the wrapper}

The commonest disguise is a single outer function wrapped around an algebraic mess. Logarithms are the classic case because the log laws are exactly a dictionary between multiplicative structure inside and additive structure outside, and additive structure is what differentiation likes. Inverse trig functions are the same phenomenon one step further along: the identity that removes the wrapper is a half-angle or double-angle formula rather than a log law, but the effect is identical.

The identities that do this work are few enough to tabulate, and tabulating them is worth half a page because recognition, not derivation, is what fails under pressure. The right-hand column is not decoration: every one of these collapses holds only on an interval, and the interval is where the traps live.

\begin{center}
\footnotesize
\setlength{\tabcolsep}{5pt}
\renewcommand{\arraystretch}{1.32}
\begin{tabular}{@{}lll@{}}
\toprule
\mh{Written as} & \mh{Collapses to} & \mh{Valid where} \\
\midrule
$\ln A+\ln B$ & $\ln(AB)$ & $A>0$, $B>0$ \\
$\ln\dfrac{A}{B}$, $\ln A^{p}$ & $\ln A-\ln B$, $p\ln A$ & $A>0$, $B>0$ \\
$\ln\bigl(\sqrt{1+x^{2}}+x\bigr)+\ln\bigl(\sqrt{1+x^{2}}-x\bigr)$ & $0$ & all $x$ \\
$\ln(\sec x+\tan x)+\ln(\sec x-\tan x)$ & $0$ & $\cos x>0$ \\
$\arctan x+\arctan(1/x)$ & $\pi/2$ & $x>0$ \\
$\arcsin x+\arccos x$ & $\pi/2$ & $-1\leq x\leq 1$ \\
$\arctan\dfrac{x}{1+\sqrt{1+x^{2}}}$ & $\tfrac12\arctan x$ & all $x$ \\
$\arcsin\dfrac{2x}{1+x^{2}}$, $\arctan\dfrac{2x}{1-x^{2}}$ & $2\arctan x$ & $\abs{x}\leq1$, $\abs{x}<1$ \\
$\sqrt{\dfrac{1-\cos x}{1+\cos x}}$ & $\tan\dfrac{x}{2}$ & $0<x<\pi$ \\
$\sin^{2}x$, $\cos^{2}x$ & $\tfrac12(1\mp\cos 2x)$ & all $x$ \\
$\sin A\cos B$ & $\tfrac12\bigl[\sin(A+B)+\sin(A-B)\bigr]$ & all $A,B$ \\
$1/\csc^{2}x$, $\sec^{2}x-\tan^{2}x$ & $\sin^{2}x$, $1$ & $\sin x\neq0$, $\cos x\neq0$ \\
$\sqrt[m]{x^{k}}$, $\ln(\cosh x+\sinh x)$ & $x^{k/m}$, $x$ & $x>0$, all $x$ \\
\bottomrule
\end{tabular}
\end{center}

Read down the second column and the pattern is uniform: a two-term expression becomes a one-term expression, a transcendental becomes an algebraic, or an expression becomes a number. That is the whole of the chapter's first half, and the rest is knowing which row you are looking at.

\tech{Log-law expansion before differentiating}{easy}
\cue A single $\ln$ or $\log$ wraps a product, a quotient, a power, or a radical: $\ln\bigl(\cos^{3}\sqrt{x}\bigr)$, $\ln\bigl(x^{x}\tan x\,\eu^{-x^{2}/\ln x}\bigr)$, and $\ln\bigl[(x^{2}+1)^{3}\sqrt{x-2}\,(x+4)^{-5}\eu^{-x}\bigr]$.
\method Apply $\ln(AB)=\ln A+\ln B$, $\ln(A/B)=\ln A-\ln B$ and $\ln(A^{p})=p\ln A$ \emph{before} differentiating, converting one chain-and-quotient monster into a sum of independent terms. Then differentiate termwise with $\dvop{x}\ln u=u'/u$. The point is that the quotient rule never appears: division inside the logarithm has already become subtraction outside it.

\begin{worked}
Take $f(x)=\ln\bigl(\cos^{3}\sqrt{x}\bigr)$ and ask for $f'(\pi^{2}/16)$. Expanding first, $f(x)=3\ln\cos\sqrt{x}$, so
\[
f'(x)=3\cdot\frac{-\sin\sqrt{x}}{\cos\sqrt{x}}\cdot\frac{1}{2\sqrt{x}}=-\frac{3\tan\sqrt{x}}{2\sqrt{x}}.
\]
At $x=\pi^{2}/16$ we have $\sqrt{x}=\pi/4$ and $\tan(\pi/4)=1$, so $f'(\pi^{2}/16)=-\dfrac{3}{2\cdot\pi/4}=-\dfrac{6}{\pi}$, numerically $-1.90986$. Note that the evaluation point was chosen so that the tangent is $1$; that is the setter telling you the expansion was intended.

A quotient case, where the saving is larger:
\[
f(x)=\ln\frac{(x^{2}+1)^{3}\sqrt{x-2}}{(x+4)^{5}\eu^{x}}=3\ln(x^{2}+1)+\tfrac12\ln(x-2)-5\ln(x+4)-x,
\]
whence $f'(x)=\dfrac{6x}{x^{2}+1}+\dfrac{1}{2(x-2)}-\dfrac{5}{x+4}-1$ and
\[
f'(3)=\frac{18}{10}+\frac12-\frac57-1=\frac{126+35-50-70}{70}=\frac{41}{70},
\]
numerically $0.585714$. Differentiating the unexpanded form means a chain rule outside a quotient rule outside two product rules, and the $\eu^{x}$ never cancels visibly.

One more, to show that the technique does not need the answer to be pretty. Take
\[
f(x)=\ln\Bigl(x^{x}\tan x\,\eu^{-x^{2}/\ln x}\Bigr)=x\ln x+\ln\tan x-\frac{x^{2}}{\ln x}.
\]
Three independent terms, each with its own one-line derivative. The first gives $\ln x+1$ by the product rule. The second gives $\dfrac{\sec^{2}x}{\tan x}=\dfrac{2}{\sin 2x}$. The third gives $\dfrac{2x\ln x-x}{(\ln x)^{2}}$ by the quotient rule, which here is applied to a genuinely irreducible two-block quotient and is therefore the right tool. So
\[
f'(x)=\ln x+1+\frac{2}{\sin 2x}-\frac{2x\ln x-x}{(\ln x)^{2}},
\]
and at $x=\eu$, where $\ln x=1$, this is $2+\dfrac{2}{\sin 2\eu}-\eu=-3.388343$. Verified against a symmetric difference quotient at step $10^{-6}$: $-3.388343$. Without the expansion, the variable exponent $x^{x}$ alone would have forced a separate logarithmic differentiation inside a chain rule inside a quotient rule.
\end{worked}

\begin{trap}
Expanding \emph{after} differentiating. A solver who writes $f'=\bigl(\text{argument}\bigr)'/\bigl(\text{argument}\bigr)$ first has committed to a quotient rule on the argument, and the resulting expression is correct but will not simplify without the very log laws that should have been used at the start. The second trap is silent: the split $\ln(A/B)=\ln A-\ln B$ requires $A>0$ and $B>0$ separately, whereas $\ln(A/B)$ needs only $A/B>0$. On the domain where both are negative you must carry absolute values, and $\ln\bigl((x-1)(x-2)\bigr)$ at $x=0$ is defined while $\ln(x-1)$ is not.
\end{trap}

\begin{seealsob}Conjugate log collapse; logarithmic differentiation of many-factor products; algebraic normalisation.\end{seealsob}

Log-law expansion assumes the sum you produce is worth having. Sometimes the sum collapses further, and in the extreme case it collapses to nothing at all. The next technique is that extreme case, and it deserves its own block because the failure to notice it is the most expensive failure in the chapter.

\tech[Conjugate log collapse and the identically constant function]{Conjugate log collapse: the function is constant}{easy}
\cue Two logarithms whose arguments are algebraic conjugates, added: $\ln\bigl(\sqrt{1+x^{2}}+x\bigr)+\ln\bigl(\sqrt{1+x^{2}}-x\bigr)$, or $\ln(\sec x+\tan x)+\ln(\sec x-\tan x)$. More generally, any expression built so that a Pythagorean or difference-of-squares identity annihilates it.
\method Combine into $\ln(AB)$ and evaluate the product before doing anything else. Conjugate pairs multiply to a difference of squares, and the whole design of these problems is that the difference of squares is $1$:
\[
\bigl(\sqrt{1+x^{2}}+x\bigr)\bigl(\sqrt{1+x^{2}}-x\bigr)=(1+x^{2})-x^{2}=1,\qquad \sec^{2}x-\tan^{2}x=1 .
\]
So $f\equiv\ln 1=0$ and $f'\equiv0$ everywhere the expression is defined.

\begin{worked}
$f(x)=\ln\bigl(\sqrt{1+x^{2}}+x\bigr)+\ln\bigl(\sqrt{1+x^{2}}-x\bigr)$; find $f'(2)$.

Combine: $f(x)=\ln\Bigl[\bigl(\sqrt{1+x^{2}}+x\bigr)\bigl(\sqrt{1+x^{2}}-x\bigr)\Bigr]=\ln 1=0$ for every real $x$. Hence $f'(2)=0$, and indeed $f'(x)=0$ identically.

It is worth seeing what the honest grind produces, because the cancellation is instructive. The first logarithm is $\operatorname{arsinh}x$, whose derivative is $1/\sqrt{1+x^{2}}$. The second is $-\operatorname{arsinh}x$, since $\sqrt{1+x^{2}}-x=1/(\sqrt{1+x^{2}}+x)$. So the two derivatives are $+1/\sqrt{1+x^{2}}$ and $-1/\sqrt{1+x^{2}}$ and they cancel. Three lines of correct algebra to reach a number that one line of identity gives for free.

The circular twin behaves the same way: $\ln(\sec x+\tan x)+\ln(\sec x-\tan x)=\ln(\sec^{2}x-\tan^{2}x)=0$, so its derivative vanishes on every interval where both factors are positive, even though each logarithm separately has the perfectly lively derivative $\pm\sec x$.
\end{worked}

\begin{trap}
Differentiating both logarithms and grinding two conjugate radical derivatives. The arithmetic is correct and the answer comes out, but a solver who does this under time pressure has spent three minutes and two sign decisions on a problem whose answer was visible. The deeper trap is the psychological one: the expression \emph{looks} expensive, so the reflex is to start computing rather than to start reading, and the design of the problem depends on that reflex.
\end{trap}

\begin{seealsob}Log-law expansion; constant scan; inverse-trig identity collapse.\end{seealsob}

Why is the identically constant function such an effective trap? Because every visual cue points the other way. The expression contains $x$ in four places; it contains a radical; it contains two transcendental functions. Nothing on the page announces constancy, and the solver's pattern-matching runs on surface complexity. The defence is not vigilance but procedure: make ``is this constant?'' the first question rather than a question you might get to, which is exactly why it sits at the top of the decision tree. And test it numerically when you can. Two evaluations at unrelated points, say $x=0.3$ and $x=1.7$, settle the matter in seconds; if the values agree, stop differentiating and start proving.

\begin{keynote}[Spot-check before you commit]
Two evaluations cost nothing and settle more than they should. If $f(0.3)$ and $f(1.7)$ agree to full precision, the function is constant and you are done. If a claimed closed form for $f'$ disagrees with a symmetric difference quotient $\bigl(f(a+h)-f(a-h)\bigr)/2h$ at $h=10^{-5}$ in the fourth digit, you have dropped a term. Every worked example in this chapter was checked that way, most by complex-step differentiation, which computes $f'(a)=\operatorname{Im}f(a+\iu h)/h$ to machine precision with no cancellation at all. The habit is not a substitute for algebra; it is what tells you which algebra to trust.
\end{keynote}

The same trap has a family. $\arctan x+\arctan(1/x)$ is $\pi/2$ for $x>0$ and $-\pi/2$ for $x<0$, so its derivative is zero on each half-line while the naive computation gives $\dfrac{1}{1+x^{2}}+\dfrac{-1/x^{2}}{1+1/x^{2}}$, which is indeed zero after simplification. Likewise $\arcsin x+\arccos x=\pi/2$, $\sin^{2}x+\cos^{2}x=1$, and $\eu^{\ln x}=x$ for $x>0$. Every one of them can be buried inside a larger expression, and every one of them collapses a subexpression to a constant that then factors out of the whole problem.

Inverse trig wrappers behave like logarithms with a different dictionary. The identity that removes them is a substitution, performed mentally rather than on paper.

\tech{Inverse-trig and half-angle identity collapse}{medium}
\cue An $\arctan$ or $\arcsin$ whose argument is a rational expression in $x$ and radicals: $\arctan\dfrac{x}{1+\sqrt{1+x^{2}}}$, $\arctan\sqrt{\dfrac{1-\cos x}{1+\cos x}}$, $\arcsin\dfrac{2x}{1+x^{2}}$.
\method Substitute $x=\tan\theta$ (for the $1+x^{2}$ shapes) or $x=\sin\theta$ (for $1-x^{2}$) in your head. The argument becomes a half-angle or double-angle expression, the inverse function cancels the direct one, and the whole wrapper collapses to a \emph{linear} function of $\arctan x$ or of $x$:
\[
\arctan\frac{x}{1+\sqrt{1+x^{2}}}=\tfrac12\arctan x,\qquad
\sqrt{\frac{1-\cos x}{1+\cos x}}=\tan\frac{x}{2}\ \ \text{on }(0,\pi),\qquad
\arcsin\frac{2x}{1+x^{2}}=2\arctan x\ \ \text{on }[-1,1].
\]
Differentiating a linear function of $\arctan x$ is then immediate.

\begin{worked}
$f(x)=\arctan\dfrac{x}{1+\sqrt{1+x^{2}}}$, find $f'(\sqrt2)$. Put $x=\tan\theta$ with $\theta\in(-\pi/2,\pi/2)$, so $\sqrt{1+x^{2}}=\sec\theta$ and
\[
\frac{x}{1+\sqrt{1+x^{2}}}=\frac{\tan\theta}{1+\sec\theta}=\frac{\sin\theta}{\cos\theta+1}=\tan\frac{\theta}{2},
\]
the last step being the standard half-angle formula. Hence $f(x)=\theta/2=\tfrac12\arctan x$ and $f'(x)=\dfrac{1}{2(1+x^{2})}$, giving $f'(\sqrt2)=\dfrac{1}{2\cdot3}=\dfrac16$.

Now the $\cos$ version. On $(0,\pi)$, $1-\cos x=2\sin^{2}(x/2)$ and $1+\cos x=2\cos^{2}(x/2)$, so the radical is $\lvert\tan(x/2)\rvert=\tan(x/2)$ there, and $g(x)=\arctan\tan(x/2)=x/2$. Therefore $g'(1)=\tfrac12$, and in fact $g'\equiv\tfrac12$ on the whole interval. Verified numerically: the symmetric difference quotient at $x=1$ with step $10^{-5}$ returns $0.5000000$.

And the double-angle version: $\arcsin\dfrac{2x}{1+x^{2}}=2\arctan x$ for $\lvert x\rvert\leq1$, so the derivative is $\dfrac{2}{1+x^{2}}$, which at $x=1/\sqrt3$ equals $\dfrac{2}{4/3}=\dfrac32$.
\end{worked}

\begin{trap}
Applying the identity on the wrong branch. Every one of these collapses is valid only on an interval, because the outer inverse function has a restricted range. Outside $(0,\pi)$ the radical $\sqrt{(1-\cos x)/(1+\cos x)}$ equals $\lvert\tan(x/2)\rvert$, so on $(\pi,2\pi)$ the function is $\pi-x/2$ and the slope is $-\tfrac12$, not $+\tfrac12$. Similarly $\arcsin\bigl(2x/(1+x^{2})\bigr)$ equals $2\arctan x$ only for $\lvert x\rvert\leq1$; for $x>1$ it is $\pi-2\arctan x$ and the derivative changes sign. Check which interval the evaluation point lies in before you quote the collapsed form.
\end{trap}

\begin{seealsob}Conjugate log collapse; trig identity reduction before high-order differentiation; log-law expansion.\end{seealsob}

\section{Identities and normal forms}

The wrapper techniques remove an outer function. The next two remove interior structure: a product of trig functions becomes a single sinusoid, and a tangle of radicals and reciprocals becomes a sum of powers. Both are about putting the expression into a normal form on which one rule, and only one rule, acts.

The trig case is where the saving is most spectacular, because it compounds with the derivative order. Differentiating a product of trig functions once is mildly annoying; differentiating it seven times is a catastrophe, unless you first notice that it was a single sine wave the whole time.

\tech{Trig identity reduction before high-order differentiation}{medium}
\cue A product or power of trigonometric functions, with a derivative of order two or more requested: $f^{(7)}$ of $\sin^{2}x$, $f''$ of $\sin 2x\cos 2x$, $f^{(5)}$ of $\sin x\cos^{3}x$.
\method Collapse to a sum of pure sinusoids using power reduction, double angle, or product to sum:
\[
\sin^{2}x=\tfrac12-\tfrac12\cos 2x,\qquad \sin 2x\cos 2x=\tfrac12\sin 4x,\qquad \sin A\cos B=\tfrac12\bigl[\sin(A+B)+\sin(A-B)\bigr].
\]
Then every order at once from the master formula
\[
\frac{\dd^{n}}{\dd x^{n}}\sin(ax+b)=a^{n}\sin\!\Bigl(ax+b+\frac{n\pi}{2}\Bigr),\qquad
\frac{\dd^{n}}{\dd x^{n}}\cos(ax+b)=a^{n}\cos\!\Bigl(ax+b+\frac{n\pi}{2}\Bigr).
\]
Each differentiation multiplies by the frequency and advances the phase by a quarter period; that is the entire content of the trigonometric derivative table.

\begin{keynote}[The quarter turn]
The four-step cycle $\sin\to\cos\to-\sin\to-\cos\to\sin$ is usually memorised as a list. It is better understood as a single statement: differentiation advances the phase by $\pi/2$ and multiplies by the frequency. That is why $\sin$ and $\cos$ need no separate table, why the cycle has period $4$, and why $f^{(n)}$ for large $n$ costs nothing more than $f'$ once the function is a sum of pure sinusoids. It also tells you which $n$ the setter will choose: $n\equiv1$ or $3\pmod 4$ to swap sine for cosine, so that a point like $\pi/4$ evaluates to something clean.
\end{keynote}

\begin{worked}
$f(x)=\sin^{2}x$; find $f^{(7)}(\pi/4)$. Power reduction gives $f(x)=\tfrac12-\tfrac12\cos 2x$. The constant contributes nothing for $n\geq1$, and
\[
f^{(7)}(x)=-\tfrac12\cdot 2^{7}\cos\!\Bigl(2x+\frac{7\pi}{2}\Bigr)=-2^{6}\cos\!\Bigl(2x-\frac{\pi}{2}\Bigr)=-64\sin 2x,
\]
using $\cos(\theta+7\pi/2)=\cos(\theta-\pi/2)=\sin\theta$. At $x=\pi/4$, $\sin(\pi/2)=1$, so $f^{(7)}(\pi/4)=-64$. Verified by contour quadrature of the seventh Cauchy coefficient: $-64.0000000$.

A harder one, to show the product-to-sum step doing real work. Take $f(x)=\sin x\cos^{3}x$ and ask for $f^{(5)}$. Write $\cos^{2}x=\tfrac12(1+\cos 2x)$ and $\sin x\cos x=\tfrac12\sin 2x$, so
\[
f(x)=\bigl(\sin x\cos x\bigr)\cos^{2}x=\tfrac12\sin 2x\cdot\tfrac12(1+\cos 2x)=\tfrac14\sin 2x+\tfrac14\sin 2x\cos 2x=\tfrac14\sin 2x+\tfrac18\sin 4x .
\]
Now the master formula, with $\sin(\theta+5\pi/2)=\cos\theta$:
\[
f^{(5)}(x)=\tfrac14\cdot2^{5}\cos 2x+\tfrac18\cdot4^{5}\cos 4x=8\cos 2x+128\cos 4x .
\]
At $x=\pi/4$ this is $8\cos(\pi/2)+128\cos\pi=0-128=-128$. Verified: $-128.0000000$. The alternative is five rounds of the product rule on a four-factor product, which generates dozens of terms that must then be collapsed by the very identities you declined to use at the start.
\end{worked}

\begin{trap}
Grinding the product rule $n$ times. The term count roughly doubles per round, so the seventh derivative of a two-factor product is a hundred-term expression before cancellation, and the cancellation is the identity you skipped. The second trap is subtler: after power reduction the constant term is invisible in every derivative, so a solver who reduces correctly but then also drops it from $f$ itself will get $f(\pi/4)$ wrong while getting every $f^{(n)}(\pi/4)$ right. Reduce for the derivatives; keep the constant for the function.
\end{trap}

\begin{seealsob}Algebraic normalisation; the cyclic $n$-th derivative pattern; log-law expansion.\end{seealsob}

Trig identities are one normal form. The other, less glamorous and far more frequently useful, is the reduction of every algebraic decoration to a rational power. This is the technique that removes the largest number of unnecessary quotient rules from a working solver's life.

\tech{Algebraic normalisation: roots to powers, quotients to sums}{easy}
\cue Radicals, nested radicals, reciprocal trig functions, or a fraction whose numerator is a \emph{sum} sitting over a single monomial denominator: $\dfrac{x^{3}+\sqrt{x}-1}{x^{2}}$, $\sqrt[3]{x^{5}}$, $\sqrt{x\sqrt{x\sqrt{x}}}$, $\dfrac{1}{\csc^{2}x}$.
\method Rewrite every root as a rational power, $\sqrt[m]{x^{k}}=x^{k/m}$, and split every sum-over-monomial quotient termwise, $\dfrac{a_{1}+\cdots+a_{n}}{g}=\dfrac{a_{1}}{g}+\cdots+\dfrac{a_{n}}{g}$ when $g$ is a single term. Send reciprocal trig functions back to base form: $1/\csc^{2}x=\sin^{2}x$, $1/\sec x=\cos x$, and use $\sec^{2}x-\tan^{2}x=1$ wherever it appears. After normalisation the power rule $\dvop{x}x^{p}=px^{p-1}$ usually does all the work, and the quotient and chain rules never enter.

\begin{worked}
Nested radicals first, because they look worst and cost least:
\[
f(x)=\sqrt{x\sqrt{x\sqrt{x}}}=\Bigl(x\cdot\bigl(x\cdot x^{1/2}\bigr)^{1/2}\Bigr)^{1/2}=\bigl(x\cdot x^{3/4}\bigr)^{1/2}=x^{7/8},
\]
so $f'(x)=\tfrac78x^{-1/8}$ and $f'(2)=\tfrac78\cdot2^{-1/8}=0.802379$. Three chain rules avoided by one exponent computation.

Now an expression engineered to look like a quotient-rule problem and to be nothing of the kind:
\[
f(x)=\frac{\bigl(\sin x+\eu^{\cos x}\bigr)\bigl(\cos^{2}x+\sin^{2}x\bigr)}{\csc^{2}x}.
\]
The middle factor is $1$ and the denominator is $\sin^{-2}x$, so $f(x)=\bigl(\sin x+\eu^{\cos x}\bigr)\sin^{2}x$: a two-factor product with no quotient anywhere. Then
\[
f'(x)=\bigl(\cos x-\sin x\,\eu^{\cos x}\bigr)\sin^{2}x+\bigl(\sin x+\eu^{\cos x}\bigr)\cdot 2\sin x\cos x .
\]
At $x=\pi/2$: $\cos x=0$, $\sin x=1$, $\eu^{\cos x}=1$. The first bracket is $0-1=-1$ and multiplies $1$; the second term carries a factor $\cos(\pi/2)=0$. So $f'(\pi/2)=-1$. Verified by complex-step differentiation: $-1.0000000$.

Finally the everyday case, a sum over a monomial:
\[
f(x)=\frac{3x^{4}-2x^{2}+5}{x^{3}}=3x-\frac{2}{x}+5x^{-3},\qquad f'(x)=3+\frac{2}{x^{2}}-15x^{-4},
\]
and $f'(1)=3+2-15=-10$. The quotient rule reaches the same place through a numerator of degree six over $x^{6}$, with a cancellation of $x^{3}$ that must be performed by hand.
\end{worked}

\begin{trap}
Splitting a quotient whose \emph{denominator} is a sum. The move $\dfrac{a+b}{g}=\dfrac ag+\dfrac bg$ is legal; the move $\dfrac{1}{a+b}=\dfrac1a+\dfrac1b$ is not, and it is the single commonest algebraic error in this family. The test is mechanical: you may split over a sum upstairs, never over a sum downstairs. A second, quieter error is losing the domain: $\sqrt{x^{2}}=\lvert x\rvert$, not $x$, so normalising $\sqrt{x^{2}}$ to $x^{1}$ silently changes the derivative from $\sgn x$ to $1$.
\end{trap}

\begin{seealsob}Log-law expansion; when the quotient rule is the wrong tool; trig identity reduction.\end{seealsob}

\section{Reading what is actually variable}

Normalisation and identity reduction share a governing idea worth naming. Both replace an expression by an equivalent one on which \emph{fewer rules} act. A nest of radicals needs three chain rules; the same function written as $x^{7/8}$ needs one power rule. A four-factor trig product needs a product rule per differentiation; the same function written as $\tfrac14\sin2x+\tfrac18\sin4x$ needs none, only a phase advance. The measure of a good normal form is exactly this: how many distinct rules survive contact with it. Aim for one.

Every technique so far assumed the whole expression depends on $x$ and asked how to reorganise it. The next assumes almost none of it does. This is a different kind of reading: not algebraic manipulation but a census of where the variable actually appears, and it is the fastest technique in the book when it applies.

\tech[Constant scan for the variable factor]{Constant scan: find the one factor that depends on $x$}{medium}
\cue A wall of intimidating symbols with at most one place where $x$ actually occurs: $\gamma$, $\eu$, $\pi$, $\Gamma$ of a number, $\cos^{1000}(\pi/\eu)$, $\sum_{n=1}^{1000}n$, $\Omega$, all arranged so the expression looks like it needs everything you know.
\method Before writing any rule, scan every sub-expression for the literal variable. Named constants, factorials and Gamma values of numbers, finite numeric sums, and transcendental functions evaluated at constants are all numbers: they differentiate to zero and factor out of the problem. Differentiate only the surviving term. If nothing survives, $f'\equiv0$.

\begin{worked}
\[
f(x)=\frac{9\,(\sin\gamma)^{99999\cos(1/x)}}{\eu^{2}+2\pi-2026\eu+\cos^{1000}(\pi/\eu)}+\sum_{n=1}^{1000}n .
\]
The census: $\gamma$ is Euler's constant, a number, so $\sin\gamma$ is a number. The entire denominator is a number; call it $K$. The finite sum $\sum_{n=1}^{1000}n=500500$ is a number. The only occurrence of $x$ in the whole display is inside $\cos(1/x)$, sitting in the exponent. So with $a=\sin\gamma$ and $E(x)=99999\cos(1/x)$ we have $f(x)=\dfrac{9}{K}a^{E(x)}+500500$ and
\[
f'(x)=\frac{9}{K}\,a^{E(x)}\ln a\cdot E'(x)=\frac{9}{K}\,(\sin\gamma)^{99999\cos(1/x)}\ln(\sin\gamma)\cdot 99999\cdot\frac{\sin(1/x)}{x^{2}},
\]
since $\dvop{x}\cos(1/x)=-\sin(1/x)\cdot(-x^{-2})=\sin(1/x)/x^{2}$. One application of the exponential rule and one chain rule, on a display that appears to need six.

A constant that is doing structural work, for contrast. In $f(x)=\cos^{1000}\bigl(x/1000\bigr)$ the exponent and the divisor are both large numbers and neither is decoration: $f'(x)=-1000\cos^{999}(x/1000)\sin(x/1000)\cdot\tfrac{1}{1000}=-\cos^{999}(x/1000)\sin(x/1000)$, and the two thousands cancel exactly. A solver who deletes them as noise loses nothing here and everything in the next problem, where they do not cancel.

A shorter instance of the same reading: $\dvop{x}\,\eu^{\sin^{2}x+\cos^{2}x}=\dvop{x}\,\eu=0$, and $\dvop{x}\bigl(x^{0}+\Gamma(7)+\arctan(\tan 1)\bigr)=0$ for $x\ne0$. Both are trivial once you notice that no exponent, no argument, and no index depends on $x$.
\end{worked}

\begin{trap}
Treating $\gamma$ or $\Omega$ as a function of $x$ because the symbol looks exotic. The reverse error is worse and rarer: assuming that a large decorative constant is inert when it is doing structural work. In the display above the exponent $99999$ multiplies the answer by $99999$; deleting it as decoration costs five orders of magnitude. Read the constants, then keep them.
\end{trap}

\begin{seealsob}Conjugate log collapse; the general power rule; algebraic normalisation.\end{seealsob}

The constant scan and the conjugate collapse are two faces of one idea: the derivative is a local statement about variation, and anything that does not vary is invisible to it. Where they differ is in how the constancy is hidden. In the constant scan it is hidden by volume: there is so much symbol on the page that the eye never audits it. In the conjugate collapse it is hidden by structure: every symbol genuinely contains $x$, and the constancy emerges only after multiplication. The first is defeated by a census, the second by an identity, and a numerical spot check defeats both.

\section{When the function is given as a process}

The last two techniques handle a different presentation entirely. So far the function has been a formula, however disguised. Now it is a recipe: a series to be summed, a product to be multiplied out, a nesting to be continued forever. In every such case the differentiation is trivial once the recipe is resolved, and essentially impossible before. Attacking these termwise is not merely slower; it is usually a different and harder problem, and it quietly assumes analytic facts that nobody has checked.

\tech[Sum the hidden series or product first]{Sum the hidden series, product, or continued fraction first}{hard}
\cue The function is presented as a $\sum$, a $\prod$, or a continued fraction, and a single derivative \emph{value} is requested. The request for one value is the tell: nobody wants a series for the answer.
\method Recognise and collapse the closed form before differentiating. The recurring engines are few. Geometric series, $\sum_{k\geq0}r^{k}=\dfrac{1}{1-r}$ for $\lvert r\rvert<1$. Doubling products, $\prod_{k=0}^{n}\bigl(1+x^{2^{k}}\bigr)=\dfrac{1-x^{2^{n+1}}}{1-x}$, which is just the binary expansion of every exponent from $0$ to $2^{n+1}-1$. The cotangent telescope, $\cot\theta-2\cot 2\theta=\tan\theta$, which sums halving tangent series. And recognising a listed Maclaurin series as an elementary function. Two further engines appear less often but are unmistakable when they do: the cyclotomic factorisation
\[
x^{2n}-1=(x^{2}-1)\prod_{k=1}^{n-1}\left(x^{2}-2x\cos\frac{k\pi}{n}+1\right),
\]
which turns a product of $n-1$ quadratics into a binomial (checked at $n=3$, $x=1.7$: both sides equal $23.137569$), and Gauss's continued fraction $\tan x=\cfrac{x}{1-\cfrac{x^{2}}{3-\cfrac{x^{2}}{5-\cdots}}}$, which is the only sane reading of a continued fraction whose partial denominators run $1,3,5,\ldots$

\begin{worked}
A geometric tail inside a radical:
\[
f(x)=\sqrt{x+1+\frac1x+\frac1{x^{2}}+\cdots}\qquad(x>1).
\]
The bracket is geometric with first term $x$ and ratio $1/x$, so it sums to $\dfrac{x}{1-1/x}=\dfrac{x^{2}}{x-1}$ and $f(x)=\dfrac{x}{\sqrt{x-1}}=x(x-1)^{-1/2}$. Then
\[
f'(x)=(x-1)^{-1/2}-\tfrac{x}{2}(x-1)^{-3/2}=(x-1)^{-3/2}\Bigl[(x-1)-\tfrac x2\Bigr]=\frac{x-2}{2(x-1)^{3/2}},
\]
so $f'(2)=0$: the function has a minimum exactly at the integer point you were asked about. Verified: $0.0000000$.

A telescoping product, done two ways as a check. Let $f(x)=(1+x)(1+x^{2})(1+x^{4})(1+x^{8})$ and find $f'(1)$. Closed form first: multiplying by $(1-x)$ telescopes, giving $f(x)=\dfrac{1-x^{16}}{1-x}=1+x+x^{2}+\cdots+x^{15}$, so $f'(1)=0+1+2+\cdots+15=120$. As a check, the logarithmic derivative of the product form at $x=1$ is $\tfrac12+\tfrac22+\tfrac42+\tfrac82=\tfrac{15}{2}$ and $f(1)=2^{4}=16$, giving $16\cdot\tfrac{15}{2}=120$ again. Verified: $120.0000000$.

Finally the tangent telescope, which is the hardest of the three and the most satisfying. Using $\tan\theta=\cot\theta-2\cot2\theta$ with $\theta=x/2^{n}$, the sum $\sum_{n\geq1}2^{-n}\tan(x/2^{n})$ telescopes to $\dfrac1x-\cot x$, since $2^{-n}\cot(x/2^{n})\to1/x$. Hence
\[
f(x)=\sum_{n\geq1}\frac{1}{2^{n}}\tan\frac{x}{2^{n}}=\frac1x-\cot x,\qquad f'(x)=-\frac{1}{x^{2}}+\csc^{2}x,
\]
and $f'(\pi/3)=-\dfrac{9}{\pi^{2}}+\dfrac{1}{3/4}=\dfrac43-\dfrac{9}{\pi^{2}}$, numerically $0.421443$. The identity itself was checked to nine digits at $x=1.1$ against $1/x-\cot x$.
\end{worked}

\begin{trap}
Differentiating termwise and then trying to sum the derivative series. This is usually a much harder sum than the original, and it silently assumes that the differentiated series converges uniformly on a neighbourhood of the point, which is exactly the hypothesis that licenses termwise differentiation. For the tangent series above, termwise differentiation produces $\sum 4^{-n}\sec^{2}(x/2^{n})$, which converges but whose closed form nobody recognises. Sum first, then differentiate.
\end{trap}

\begin{seealsob}Self-similar fixed-point resolution; algebraic normalisation; logarithmic differentiation of many-factor products.\end{seealsob}

\begin{keynote}[Process forms are not harder, they are earlier]
A function given as a series, a product, or an infinite nesting is not a harder differentiation problem. It is a problem that has not been stated yet. The statement is the closed form, and finding it is algebra, summation, or fixed-point reasoning, none of which is calculus. Once it is found the differentiation is routine, which is the tell: if the derivative you are computing is itself hard, you have almost certainly not finished the previous step. This is also why these problems ask for a single \emph{value}. A value can be checked; a series cannot.
\end{keynote}

Summing a series requires a formula for the partial sums. Infinite nestings do not have partial sums in any useful sense, but they have something better: they contain themselves, and that self-reference converts an infinite object into a finite equation in one line.

\tech{Self-similar fixed-point resolution}{hard}
\cue An infinite nesting that repeats itself: a power tower $x^{x^{x^{\cdots}}}$, a nested radical $\sqrt{x+\sqrt{x+\cdots}}$, an infinite logarithmic tower $\ln\bigl(x\ln(x\ln\cdots)\bigr)$, or the limit of an iterated map.
\method Name the whole object $t$ and exploit self-similarity: the thing inside the outermost operation is the whole object again. That gives a finite equation $t=\Phi(x,t)$. Solve it if you can; otherwise differentiate it implicitly, which is almost always enough. For the tower, $t=x^{t}$; for the log tower, $t=\ln(xt)$, that is $\eu^{t}=xt$; for an iterated map with limit $L$, solve $L=\varphi(L)$ and note that $L$ is then a genuine constant and the calculus is over.

\begin{worked}
The log tower, where the inverse function is the clean object. Let $f(x)=\ln\bigl(x\ln(x\ln(\cdots))\bigr)$. Self-similarity gives $f=\ln\bigl(xf\bigr)$, hence $\eu^{f}=xf$. Now solve for $x$ rather than for $f$: if $y=f(x)$ then $x=\dfrac{\eu^{y}}{y}$, so the inverse function is explicit,
\[
f^{-1}(y)=\frac{\eu^{y}}{y},\qquad \bigl(f^{-1}\bigr)'(y)=\frac{\eu^{y}y-\eu^{y}}{y^{2}}=\frac{\eu^{y}(y-1)}{y^{2}},
\]
and at $y=\tfrac32$ this is $\dfrac{\eu^{3/2}\cdot\tfrac12}{\tfrac94}=\dfrac{2\eu^{3/2}}{9}$, numerically $0.995931$. Verified by complex-step differentiation of $\eu^{y}/y$.

The power tower, done implicitly. Let $t=x^{x^{x^{\cdots}}}$, so $t=x^{t}$ and $\ln t=t\ln x$. Differentiating both sides in $x$,
\[
\frac{t'}{t}=t'\ln x+\frac{t}{x}\quad\Longrightarrow\quad t'\Bigl(\frac1t-\ln x\Bigr)=\frac tx\quad\Longrightarrow\quad t'=\frac{t^{2}}{x\bigl(1-t\ln x\bigr)} .
\]
At $x=\sqrt2$ the tower converges to $t=2$, since $(\sqrt2)^{2}=2$, and $t\ln x=2\cdot\tfrac12\ln2=\ln 2$. So $t'(\sqrt2)=\dfrac{4}{\sqrt2\,(1-\ln2)}=9.21754$. Verified by iterating the tower four thousand times and taking a symmetric difference quotient with step $10^{-6}$: $9.21754$.


The nested radical, which is the gentlest member of the family and the best place to see the mechanism. Let $t=\sqrt{x+\sqrt{x+\sqrt{x+\cdots}}}$. Everything under the outermost radical is $x$ plus the whole object again, so $t^{2}=x+t$, a quadratic. Taking the root that is positive for $x>0$,
\[
t=\frac{1+\sqrt{1+4x}}{2},\qquad \frac{\dd t}{\dd x}=\frac{1}{\sqrt{1+4x}} ,
\]
or, without solving at all, differentiate $t^{2}=x+t$ implicitly to get $2tt'=1+t'$, hence $t'=\dfrac{1}{2t-1}$, which is the same thing. At $x=2$ the fixed point is $t=2$ and $t'=\tfrac13$. Verified by iterating the radical two thousand times and taking a symmetric difference quotient: $0.3333333$. The implicit route is the one to prefer, because it works when the fixed-point equation cannot be solved in closed form.
\end{worked}

\begin{trap}
Forgetting to ask which fixed point the construction converges to, and whether it converges at all. The tower $x^{x^{x^{\cdots}}}$ converges only for $\eu^{-\eu}\leq x\leq \eu^{1/\eu}$, and the equation $t=x^{t}$ has two roots on part of that range; at $x=\sqrt2$ the roots are $t=2$ and $t=4$, and the tower converges to $2$. Writing $t'$ in terms of the wrong root gives a confident wrong number. The second trap is the reverse of the usual one: after $x$ is replaced by a specific number the nesting may collapse to a constant, so the calculus becomes trivial rather than harder, and a solver braced for difficulty can walk straight past that.
\end{trap}

\begin{seealsob}Sum the hidden series, product, or continued fraction first; power towers and iterated exponentials; conjugate log collapse.\end{seealsob}

\subsection{The scan applied end to end}

Every technique above has been shown in isolation. Real problems stack them, and the stacking is the point: a setter who wants a hard problem applies three identities backwards, not one. Here is a single expression that needs four of the chapter's moves in sequence, and which takes two lines once they are made and rather more than two pages if they are not. For $x>0$, let
\[
f(x)=\ln\!\left(\frac{\bigl(\sin 2x\cos 2x\bigr)^{3}\ \sqrt[4]{x^{6}}}{\eu^{\arctan x+\arctan(1/x)}}\right).
\]

\emph{Move one, log laws.} The outer logarithm wraps a product, a quotient, and two powers, so expand:
\[
f(x)=3\ln\bigl(\sin 2x\cos 2x\bigr)+\tfrac{6}{4}\ln x-\bigl(\arctan x+\arctan(1/x)\bigr),
\]
the exponential in the denominator having been annihilated by $\ln\eu^{u}=u$.

\emph{Move two, the constant scan.} The last bracket is $\pi/2$ for every $x>0$. It contributes nothing to $f'$ and can be struck out now.

\emph{Move three, a trig identity.} $\sin 2x\cos 2x=\tfrac12\sin 4x$, so the first term is $3\ln\tfrac12+3\ln\sin 4x$, and the constant $3\ln\tfrac12$ also dies.

\emph{Move four, radicals to powers.} $\sqrt[4]{x^{6}}=x^{3/2}$, so the middle term is $\tfrac32\ln x$, which was already the case but is now visibly so.

What remains is $f(x)=3\ln\sin 4x+\tfrac32\ln x+\text{const}$, and one differentiation finishes it:
\[
f'(x)=3\cdot\frac{4\cos 4x}{\sin 4x}+\frac{3}{2x}=12\cot 4x+\frac{3}{2x}.
\]
At $x=\pi/16$ we have $4x=\pi/4$ and $\cot(\pi/4)=1$, so
\[
f'\!\left(\frac{\pi}{16}\right)=12+\frac{3}{2\cdot\pi/16}=12+\frac{24}{\pi}=19.639437 .
\]
Verified by complex-step differentiation: $19.639437$. Count what was avoided: a chain rule around a quotient rule around two product rules around a chain rule, and an $\arctan$ pair whose derivatives cancel only after being written out. The evaluation point $\pi/16$ was, as always, the setter's signature: it is the point at which the surviving cotangent is $1$.

One closing remark on procedure, since the chapter has been a list of moves and the moves are worth nothing without the habit that triggers them. The habit is cheap to install: after reading the problem and before writing anything, say out loud what the expression \emph{is}. ``It is a logarithm of a quotient.'' ``It is a product of two sinusoids.'' ``It is a constant.'' The sentence forces the census that the decision tree formalises, and in the overwhelming majority of cases it also names the technique. When the sentence comes out as ``it is a genuinely irreducible composition of three functions'', you have learned something too, and the next chapter is about what to do then.

The chapter's claim, stated once more in its sharpest form: there is no such thing as a hard derivative, only a hard expression. Every technique here is a way of making the expression easier before the derivative is taken, and every one of them is an identity you already know, applied in the direction the setter did not. What the chapter adds is the order in which to try them and the discipline to try them at all. The next chapter takes up the residue, the expressions that genuinely do not collapse, and asks a different question: given that you must differentiate, which rule system should you use?

\section{Recognition matrix}
\begin{recog}{@{}p{0.30\textwidth}p{0.36\textwidth}p{0.28\textwidth}@{}}
\rowcolor{mist}\mh{If you see} & \mh{Reach for} & \mh{If that fails} \\
\midrule
\endhead
$\ln$ of a product, quotient, power or radical & Log laws first, then differentiate termwise & Check positivity of each split factor \\
Two logs with conjugate arguments, added & Combine: the product is usually $1$, so $f'\equiv0$ & Evaluate $f$ at two points and compare \\
$\arctan$ or $\arcsin$ of a rational-in-radicals argument & $x=\tan\theta$ or $x=\sin\theta$; a half-angle collapse & Identify the interval; the branch may flip the sign \\
$f^{(n)}$ of a product or power of trig functions & Power reduction to $\sum c_{k}\sin(a_{k}x+b_{k})$ & Master formula: $a^{n}\sin(ax+b+n\pi/2)$ \\
Radicals or nested radicals of powers of $x$ & Collapse the exponents; then the power rule alone & Watch $\sqrt{x^{2}}=\abs{x}$ \\
A sum over a single monomial denominator & Split termwise; never split over a sum downstairs & Write $uv^{-1}$ and use product $+$ chain \\
$1/\csc^{2}x$, $\sec^{2}x-\tan^{2}x$, $\cos^{2}x+\sin^{2}x$ & Reciprocal trig to base form; Pythagorean identities & Substitute a numeric $x$ to confirm the value \\
A wall of $\gamma$, $\pi$, $\Gamma(7)$, $\sum_{n=1}^{1000}n$ & Census for the literal $x$; everything else is a number & If no $x$ survives, $f'\equiv0$ \\
$\arctan x+\arctan(1/x)$, $\arcsin x+\arccos x$ & A constant on each interval; derivative $0$ & Confirm the constant's value per branch \\
$f$ given as $\sum$ or $\prod$, one value wanted & Sum or telescope to closed form, then differentiate & Log-differentiate the finite product instead \\
$\prod_{k=0}^{n}(1+x^{2^{k}})$ & Multiply by $(1-x)$: it telescopes & Logarithmic derivative of the product \\
$\sum 2^{-n}\tan(x/2^{n})$ and relatives & Cotangent telescope $\tan\theta=\cot\theta-2\cot2\theta$ & Check the identity numerically first \\
$x^{x^{x^{\cdots}}}$ or $\sqrt{x+\sqrt{x+\cdots}}$ & Set $t=\Phi(x,t)$, differentiate implicitly & Solve for $x$ in terms of $t$ and invert \\
An infinite nesting at a specific $x$ & Check convergence and which fixed point & The nesting may be a constant there \\
Nothing collapses after the four questions & Go to the rule systems: chain, product, quotient & Logarithmic differentiation as a solvent \\
\end{recog}
